\section{从实数轴到函数}
\label{sec:02-Calculus/01-Functions-and-Precalculus/01}

一份模型接口文档写着：``输出一个 \(0\) 到 \(1\) 之间的置信分数。''

上线三个月后有人抽了十万条预测出来看，最大 \(0.87\)，最小 \(0.02\)，两头空着的区间一次都没出现。文档没写错。它承诺的是输出\emph{可以}落在哪里，而下游把``分数超过 \(0.9\) 自动放行''写成了规则，于是这条规则从来没有触发过。一个月的工期，浪费在一个语义混淆上。

\subsection{陪域是承诺，值域是事实}

数学早就把这两件事分开命名了。写 \(f:D\to C\) 时，\(D\) 是定义域，是允许喂进去的输入集合；\(C\) 是陪域，是你事先声明输出不会跑出去的范围。而值域 \(f(D)\) 是所有真正被打到的输出，它通常只是陪域的一部分。\(f(x)=x^2\) 作为 \(f:\R\to\R\)，陪域是全体实数，值域却只有 \([0,\infty)\)。

对应规则本身只有两条要求：\(D\) 里每一个 \(x\) 都得有输出，且这个输出唯一。前一条禁止``这个输入不知道该给什么''，后一条禁止同一个输入既给 \(3\) 又给 \(5\)。

\begin{center}
\begin{tikzpicture}[x=1cm,y=1cm,font=\small,>=stealth]
  \draw[black!45] (0,0) ellipse (1.3 and 2.0);
  \node[above,font=\footnotesize] at (0,2.1) {定义域 \(D\)};
  \draw[black!45] (5.6,0) ellipse (1.5 and 2.4);
  \node[above,font=\footnotesize] at (5.6,2.5) {陪域 \(C\)};
  \draw[black!55,fill=blue!12] (5.55,-0.55) ellipse (1.05 and 1.45);
  \node[font=\footnotesize] at (5.55,-0.55) {值域 \(f(D)\)};
  \fill (-0.15,1.15) circle (1.6pt);
  \fill (-0.15,0.05) circle (1.6pt);
  \fill (-0.15,-1.15) circle (1.6pt);
  \fill (5.6,0.35) circle (1.6pt);
  \fill (5.25,-1.15) circle (1.6pt);
  \draw[->,black!70] (0.05,1.15) -- (5.45,0.42);
  \draw[->,black!70] (0.05,0.05) -- (5.45,0.30);
  \draw[->,black!70] (0.05,-1.15) -- (5.10,-1.15);
  \fill[red!70!black] (5.7,1.75) circle (1.6pt);
  \node[right,font=\footnotesize,red!55!black] at (5.85,1.75) {没有原像};
\end{tikzpicture}

\emph{阴影块是输出实际覆盖的地方，外面那圈只是承诺；红点在陪域内却没有任何输入能打到它。另外注意上面两支箭头落在同一个点上——一旦发生这种碰撞，光看输出就再也认不出输入是谁了。}
\end{center}

这张图之后还会用两次。碰撞那件事是可逆性的全部障碍，本章第三节会正面处理；红点那件事在评估模型时天天发生——校准曲线画到 \(0.9\) 以上没有数据点，不是模型在那一段表现差，而是它根本没去过那里。

\subsection{定义域是规则的一部分}

同一个代数表达式配上不同的定义域，是不同的函数。\(f(x)=\sqrt{1-x^2}\) 在实数范围内自动只对 \([-1,1]\) 有定义；若问题只关心右半边，你可以主动把定义域缩到 \([0,1]\)，此时图像、值域、能讨论的趋近方向全都随之改变。缩定义域不是数学上偷懒，而是在声明``在我的问题里只有这段输入有意义''。

也不要以为函数必须有公式。每天的气温记录是函数，输入时刻输出温度；算法运行时间关于输入规模是函数；``前 \(3\) 公里 \(10\) 元，之后每公里 \(2\) 元''是函数；由微分方程定义、写不出闭式解的轨迹同样是函数。闭式表达式（用有限个初等运算写出来的式子）是一种方便，不是资格证。正态分布的累积分布函数就没有初等闭式，这丝毫不妨碍统计学天天用它。

\subsection{距离把``接近''变成可以检查的量}

微积分从头到尾都在说``\(x\) 靠近 \(a\) 时会怎样''。顺序关系不够用：\(0.5<1\) 和 \(0.999<1\) 在顺序看来是同一件事，可我们显然想区分这两个位置。需要的是距离，而实数轴上的距离就是 \(|x-a|\)。

于是``\(x\) 落在 \(a\) 附近''有了精确写法，并且它和区间记号是同一件事的两种说法：

\[
|x-a|<r
\quad\Longleftrightarrow\quad
a-r<x<a+r
\quad\Longleftrightarrow\quad
x\in(a-r,\,a+r).
\]

区间的方括号表示端点算数，圆括号表示不算，无穷一侧永远用圆括号——无穷不是实数，只是方向标记，没有``包含''可言。半开半闭在分段函数的边界上格外要紧，后面马上会看到。

这条等价是整个极限定义的几何核心。为了让 \(f(x)\) 落在目标值 \(L\) 周围的误差窗里，你要找到 \(x\) 周围一个足够小的窗；两个窗都由绝对值度量，同一把尺子在输入端和输出端各用一次。

\subsection{实数轴没有缺口，这件事后面处处在用}

画数轴时你默认它没有洞，但有理数轴确实有洞：单位正方形的对角线长 \(\sqrt2\)，这个位置在有理数里空着。实数补上了所有这类空位，这条性质叫完备性，最常用的表述是任何有上界的非空实数集必有最小上界。

现在不需要证明它（\hyperref[ch:097]{``实数为什么需要完备''}一章会正式展开），但值得知道它在替谁站岗。介值定理说连续函数必须穿过中间的每一个值，最值定理说闭区间上的连续函数一定取到最大最小值——这两条``看起来显然''的定理，证明最后都落回完备性。纸上的线看着连续只是错觉，屏幕的像素也是有限的；能担保这件事的从来不是图，而是完备性。

\subsection{图像是窗户，不是地基}

实函数的图像是点集 \(\{(x,f(x)):x\in D\}\)。竖线检验直接来自``输出唯一''这条要求：竖线穿过图像两次，就意味着同一个 \(x\) 配了两个 \(y\)。

图像的用处很实在：升降、拐点、渐近行为一眼可见。危险也在这里，看见了就不再追问凭什么。有三类欺骗值得记住。分辨率会藏细节：\(\sin(1/x)\) 在 \(0\) 附近振荡频率无限增大，任何屏幕只能画出有限次振荡，你看到的``密集波浪''和数学上的行为是两回事。缺口会消失：\((x^2-1)/(x-1)\) 在 \(x\neq1\) 时等于 \(x+1\)，画出来是一条直线，那个洞小到肉眼看不见。尺度会误导：\(x^{100}\) 和 \(x^2\) 在 \([-1,1]\) 上都像平底碗，可在 \(x=0.99\) 处一个取 \(0.366\)，另一个取 \(0.980\)，边界附近的行为天差地别。

用法因此是固定的：图像用来生成猜想，代数和极限用来确认或推翻。

\subsection{分界点是分段函数里唯一值得盯的位置}

有些依赖关系写不成一个统一公式，但完全符合函数的定义。绝对值就是最短的例子，\(x\ge0\) 时取 \(x\)，\(x<0\) 时取 \(-x\)：同一个函数，两条规则，每个输入仍然只有一个输出。

看分段函数时，信息几乎全部集中在规则切换的那个点上。取

\[
f(x)=
\begin{cases}
x+1, & x<0,\\[4pt]
x^2, & x\ge 0.
\end{cases}
\]

从左边靠近 \(0\)，输出依次是 \(0.5,\,0.9,\,0.99\)，往 \(1\) 去；从右边靠近，输出是 \(0.25,\,0.01,\,0.0001\)，往 \(0\) 去；而 \(f(0)\) 按规则等于 \(0\)。左边飘向一个值，右边和函数值本身停在另一个值，图像在这里裂开一道口子。含绝对值的表达式是分段函数的另一个常见来源，找分界点只需要问表达式在哪里变号。

请留意刚才那段描述里悄悄用掉的东西：``从左边靠近''和``从右边靠近''。绝对值只量距离，不分方向，\(|x-0|\) 对 \(-0.01\) 和 \(0.01\) 给出同一个数，可这两侧的函数值去向完全不同。要把这道裂口说清楚，需要给趋近加上方向，还需要一种不依赖``到底取不取得到''的语言。

在此之前，先把手上的函数变多一点。下一节从少数几个熟悉的图像出发，用平移、缩放和串联生成一整族新函数，并说明为什么横向的加号会把图像往左推。
